\SourceLeaf{25}
\par
The attraction in the direction parallel and opposite to the $y$-coordinate axis is obtained by replacing $a$ with $b$ in these formulas. It is clear that these two forces are equivalent to a single force whose direction is normal to the axis $2C$ and whose magnitude is found by replacing $a$ in the formula just given with the distance of the attracted point from that axis.

Finally, the attraction in the direction parallel and opposite to the $z$-coordinate axis, that is, normal to the equator, is, in the case where $B=A$,
\[
=\frac{4\pi cAA}{CC}\int\frac{tt\mathrm{d}t}{1-(1-\frac{AA}{CC})tt}
\]
from which, if $C<A$, setting $\frac{C}{A}=\cos\varphi$ as above gives
\[
=\frac{4\pi c\cos\varphi}{\sin\varphi^3}(\operatorname{tang}\varphi-\varphi)
\]
whereas if $C>A$ we obtain
\[
\frac{4\pi cAAC}{(CC-AA)^{\frac32}}\log\frac{C+\BareRoot{(CC-AA)}}{A}-\frac{4\pi cAA}{CC-AA}
\]
Lastly, if all three $A,B,C$ are equal, that is, if the body is a sphere, the attractions along the three principal directions become
\[
\tfrac43\pi a,\quad\tfrac43\pi b,\quad\tfrac43\pi c
\]
That is, they are identical to the attractions that a spherical core with the attracted point on its surface would exert if its mass were concentrated at its center. It follows at once that \textit{exterior} points are attracted by a sphere just as if its entire mass were at its center, as \textsc{Newton} first showed.

\par\bigskip\Needspace{6\baselineskip}
\pdfbookmark[1]{Addendum}{attractionis-additamentum}
\centerline{ADDENDUM.}\medskip
\par
After this had already been written, the distinguished \textsc{Laplace} drew attention to an excellent paper by the eminent \textsc{Ivory} in the \textit{Philosophical Transactions} for 1809. There the same subject is treated by a method entirely different from those used by the distinguished \textsc{Laplace} and \textsc{Legendre}. With the greatest elegance, that geometer showed how to reduce the attraction on an exterior point to the attraction on an interior point, that is, how to reduce the part of the problem always regarded as harder to the easier part. But the method by which
